\section{研究背景}

\begin{frame}{研究背景：双碳目标需要金融体系提供持续支持}
  \begin{columns}[T,onlytextwidth]
    \column{0.56\textwidth}
    \begin{itemize}
      \item 气候变化治理已从环境议题上升为 \hlb{经济结构转型议题}。
      \item 中国提出 \hl{2030 碳达峰、2060 碳中和}，低碳转型面临巨大资金缺口。
      \item 绿色金融通过信贷、债券、基金、保险等工具，引导资本流向绿色产业。
      \item 关键问题不只是“规模有没有增长”，而是 \hlt{是否真的带来了减排效果}。
    \end{itemize}

    \begin{block}{本文关注}
      绿色金融发展是否能够显著降低各省 \hlb{碳排放强度}，
      以及这种效果在东中西部地区是否存在明显差异。
    \end{block}

    \column{0.4\textwidth}
    \begin{exampleblock}{现实动机}
      \begin{itemize}
        \item 地区产业结构差异大
        \item 金融市场发育程度不同
        \item 政策资源投放效果可能不均衡
      \end{itemize}
    \end{exampleblock}
  \end{columns}
\end{frame}

\begin{frame}{研究意义：回答“绿色金融是否有效、何处更有效”}
  \begin{enumerate}
    \item 从学术上，为 \hlb{绿色金融与环境经济学交叉研究} 提供省级面板实证证据。
    \item 从政策上，检验绿色金融发展是否真正转化为 \hl{碳减排绩效}。
    \item 从实践上，识别东中西部效果差异，为 \hlt{差异化政策设计} 提供依据。
  \end{enumerate}

  \vspace{0.6em}
  \begin{alertblock}{答辩主线}
    政策背景提出问题，理论机制建立假设，实证结果检验效果，最后回到政策工具优化。
  \end{alertblock}
\end{frame}
